\documentclass[12pt, ukrainian]{article}
\usepackage[a4paper]{geometry}
\setlength\parindent{0mm}
%\setlength\parskip{4mm}
\geometry{top=1.3cm,bottom=1.3cm,left=1.5cm,right=1.5cm}
\pagestyle{empty}
\usepackage[T2A]{fontenc}
\usepackage[utf8]{inputenc}
\usepackage[ukrainian]{babel}
\begin{document}
%begin
{{\begin {scriptsize}\par
{ \center  Київський національний університет імені Тараса Шевченка \par}
{\parbox {5 cm}{Освітня програма \hfill}{\parbox {7 cm}{Інформатика\hfill}}\parbox {6 cm}{\hfill Семестр 1 (маг)}}\par
{\parbox {5 cm} {Навчальний предмет \hfill}\parbox {7 cm}{Інформаційні мережі \hfill}\parbox {6cm}{\hfill}\par}\vspace{-15pt} \end{scriptsize}}{\center Екзаменаційний білет \No \textbf { 1 }\par}}
{
\qquad\textbf{Дати розгорнуту та вичерпну відповідь на питання}, висвітливши необхідні означення, властивості та за необхідності провівши порівняльний аналіз.

%beginitem
1. Методи забезпечення належного рівня QoS. Методи боротьби з перенавантаженням.

%enditem


%\vspace{2mm}
\qquad\textbf{Визначити істинність таких тверджень}.

%beginitem
2. Модель OSI стосується мережі комутації каналів.

%enditem

%beginitem
3. На фізичному, канальному та мережевому рівнях об'єкти взаємодіють через інтерфейси, а на транспортному та прикладному рівнях через протоколи.

%enditem

%beginitem
4. При передачі мультимедіа в потоковому режимі за втрати пакетів зазвичай виконується повторне відправлення втрачених пакетів.

%enditem

%beginitem
5. Велика варіація затримки є критичною для передачі мультимедіа в режимі реального часу.

%enditem

%beginitem
6. Використання ChatGPT на іспиті є неавторизованим доступом до інформації.

%enditem
}
{\smallskip\smallskip \begin {scriptsize} Затверджено на засіданні кафедри теоретичної кібернетики, протокол \No 2  від   07.10.2024 р.\par\smallskip\smallskip\smallskip\smallskip
{\parbox {6.5 cm} {Зав. кафедри \dotfill Ю.В.~Крак}}\hspace {0.7cm} {\hbox to 6.5 cm { Екзаменатор \dotfill Т.О.~Карнаух}} \end{scriptsize}}
%end
\vspace{0.9cm}
%begin
{{\begin {scriptsize}\par
{ \center  Київський національний університет імені Тараса Шевченка \par}
{\parbox {5 cm}{Освітня програма \hfill}{\parbox {7 cm}{Інформатика\hfill}}\parbox {6 cm}{\hfill Семестр 1 (маг)}}\par
{\parbox {5 cm} {Навчальний предмет \hfill}\parbox {7 cm}{Інформаційні мережі \hfill}\parbox {6cm}{\hfill}\par}\vspace{-15pt} \end{scriptsize}}{\center Екзаменаційний білет \No \textbf { 2 }\par}}
{
\qquad\textbf{Дати розгорнуту та вичерпну відповідь на питання}, висвітливши необхідні означення, властивості та за необхідності провівши порівняльний аналіз.

%beginitem
1. Стек протоколів TCP/IP. Співвідношення з моделлю OSI.

%enditem


%\vspace{2mm}
\qquad\textbf{Визначити істинність таких тверджень}.

%beginitem
2. Маршрутизатор зазвичай працює на прикладному рівні.

%enditem

%beginitem
3. Мережевий протокол є фізичним інтерфейсом між об'єктами одного рівня.

%enditem

%beginitem
4. Протокол UDP працює без встановлення з'єднання.

%enditem

%beginitem
5. Еластичний трафік є чутливим до затримок.

%enditem

%beginitem
6. Використання ChatGPT на іспиті є авторизованим доступом до інформації.

%enditem
}
{\smallskip\smallskip \begin {scriptsize} Затверджено на засіданні кафедри теоретичної кібернетики, протокол \No 2  від   07.10.2024 р.\par\smallskip\smallskip\smallskip\smallskip
{\parbox {6.5 cm} {Зав. кафедри \dotfill Ю.В.~Крак}}\hspace {0.7cm} {\hbox to 6.5 cm { Екзаменатор \dotfill Т.О.~Карнаух}} \end{scriptsize}}
%end
\vspace{0.9cm}
%begin
{{\begin {scriptsize}\par
{ \center  Київський національний університет імені Тараса Шевченка \par}
{\parbox {5 cm}{Освітня програма \hfill}{\parbox {7 cm}{Інформатика\hfill}}\parbox {6 cm}{\hfill Семестр 1 (маг)}}\par
{\parbox {5 cm} {Навчальний предмет \hfill}\parbox {7 cm}{Інформаційні мережі \hfill}\parbox {6cm}{\hfill}\par}\vspace{-15pt} \end{scriptsize}}{\center Екзаменаційний білет \No \textbf { 3 }\par}}
{
\qquad\textbf{Дати розгорнуту та вичерпну відповідь на питання}, висвітливши необхідні означення, властивості та за необхідності провівши порівняльний аналіз.

%beginitem
1. Характеристики інфокомунікаційних мереж.

%enditem


%\vspace{2mm}
\qquad\textbf{Визначити істинність таких тверджень}.

%beginitem
2. Маршрутизатор зазвичай працює на мережевому рівні.

%enditem

%beginitem
3. Мережевий протокол містить інформацію про кількість даних, передану мережею за певну одиницю часу.

%enditem

%beginitem
4. Протокол HTTP має механізм контролю перенавантаження (congestion).

%enditem

%beginitem
5. За інтенсивності трафіка 2 мережа ще може нормально функціонувати, а за інтенсивності 5 та більше варто вживати термінових заходів для підтримки працездатності мережі.

%enditem

%beginitem
6. Співробітник ректорату на іспиті вирішив перевірити документи студентів. Це він авторизує студентів.

%enditem
}
{\smallskip\smallskip \begin {scriptsize} Затверджено на засіданні кафедри теоретичної кібернетики, протокол \No 2  від   07.10.2024 р.\par\smallskip\smallskip\smallskip\smallskip
{\parbox {6.5 cm} {Зав. кафедри \dotfill Ю.В.~Крак}}\hspace {0.7cm} {\hbox to 6.5 cm { Екзаменатор \dotfill Т.О.~Карнаух}} \end{scriptsize}}
%end
\newpage
%begin
{{\begin {scriptsize}\par
{ \center  Київський національний університет імені Тараса Шевченка \par}
{\parbox {5 cm}{Освітня програма \hfill}{\parbox {7 cm}{Інформатика\hfill}}\parbox {6 cm}{\hfill Семестр 1 (маг)}}\par
{\parbox {5 cm} {Навчальний предмет \hfill}\parbox {7 cm}{Інформаційні мережі \hfill}\parbox {6cm}{\hfill}\par}\vspace{-15pt} \end{scriptsize}}{\center Екзаменаційний білет \No \textbf { 4 }\par}}
{
\qquad\textbf{Дати розгорнуту та вичерпну відповідь на питання}, висвітливши необхідні означення, властивості та за необхідності провівши порівняльний аналіз.

%beginitem
1. Маршрутизація, її задачі. Алгоритми та протоколи маршрутизації.

%enditem


%\vspace{2mm}
\qquad\textbf{Визначити істинність таких тверджень}.

%beginitem
2. Маршрутизатор зазвичай працює на канальному (data link) рівні.

%enditem

%beginitem
3. Мережевий протокол містить інформацію про відхилення параметрів мережі від заданих.

%enditem

%beginitem
4. Bluetooth є протоколом канального рівня.

%enditem

%beginitem
5. Для передачі документів мережею велика варіація затримки некритична.

%enditem

%beginitem
6. Співробітник ректорату на іспиті вирішив перевірити документи студентів. Це він автентифікує студентів.

%enditem
}
{\smallskip\smallskip \begin {scriptsize} Затверджено на засіданні кафедри теоретичної кібернетики, протокол \No 2  від   07.10.2024 р.\par\smallskip\smallskip\smallskip\smallskip
{\parbox {6.5 cm} {Зав. кафедри \dotfill Ю.В.~Крак}}\hspace {0.7cm} {\hbox to 6.5 cm { Екзаменатор \dotfill Т.О.~Карнаух}} \end{scriptsize}}
%end
\vspace{0.9cm}
%begin
{{\begin {scriptsize}\par
{ \center  Київський національний університет імені Тараса Шевченка \par}
{\parbox {5 cm}{Освітня програма \hfill}{\parbox {7 cm}{Інформатика\hfill}}\parbox {6 cm}{\hfill Семестр 1 (маг)}}\par
{\parbox {5 cm} {Навчальний предмет \hfill}\parbox {7 cm}{Інформаційні мережі \hfill}\parbox {6cm}{\hfill}\par}\vspace{-15pt} \end{scriptsize}}{\center Екзаменаційний білет \No \textbf { 5 }\par}}
{
\qquad\textbf{Дати розгорнуту та вичерпну відповідь на питання}, висвітливши необхідні означення, властивості та за необхідності провівши порівняльний аналіз.

%beginitem
1. Модель OSI. Рівні та їх задачі. Принципи будови моделі.

%enditem


%\vspace{2mm}
\qquad\textbf{Визначити істинність таких тверджень}.

%beginitem
2. Маршрутизатор зазвичай працює на транспортному рівні.

%enditem

%beginitem
3. Мережевий протокол є логічним інтерфейсом між об'єктами одного рівня.

%enditem

%beginitem
4. У протоколі HTTP встановлюється з'єднання між комп'ютерами користувачів.

%enditem

%beginitem
5. Пересічний користувач під час веб-сьорфінгу породжує пульсуючий трафік.

%enditem

%beginitem
6. Співробітник ректорату на іспиті вирішив перевірити документи студентів. Це він ідентифікує студентів.

%enditem
}
{\smallskip\smallskip \begin {scriptsize} Затверджено на засіданні кафедри теоретичної кібернетики, протокол \No 2  від   07.10.2024 р.\par\smallskip\smallskip\smallskip\smallskip
{\parbox {6.5 cm} {Зав. кафедри \dotfill Ю.В.~Крак}}\hspace {0.7cm} {\hbox to 6.5 cm { Екзаменатор \dotfill Т.О.~Карнаух}} \end{scriptsize}}
%end
\vspace{0.9cm}
%begin
{{\begin {scriptsize}\par
{ \center  Київський національний університет імені Тараса Шевченка \par}
{\parbox {5 cm}{Освітня програма \hfill}{\parbox {7 cm}{Інформатика\hfill}}\parbox {6 cm}{\hfill Семестр 1 (маг)}}\par
{\parbox {5 cm} {Навчальний предмет \hfill}\parbox {7 cm}{Інформаційні мережі \hfill}\parbox {6cm}{\hfill}\par}\vspace{-15pt} \end{scriptsize}}{\center Екзаменаційний білет \No \textbf { 6 }\par}}
{
\qquad\textbf{Дати розгорнуту та вичерпну відповідь на питання}, висвітливши необхідні означення, властивості та за необхідності провівши порівняльний аналіз.

%beginitem
1. QoS. Основні типи QoS. Угода про якість обслуговування (SLA).

%enditem


%\vspace{2mm}
\qquad\textbf{Визначити істинність таких тверджень}.

%beginitem
2. Модель OSI стосується мережі комутації пакетів.

%enditem

%beginitem
3. Мережевий протокол є логічним інтерфейсом між різними рівнями стеку протоколів.

%enditem

%beginitem
4. FTP є протоколом транспортного рівня.

%enditem

%beginitem
5. Передача документів породжує еластичний трафік.

%enditem

%beginitem
6. В аудиторію зайшов співробітник ректорату і представився. Це він виконав аутентифікацію.

%enditem
}
{\smallskip\smallskip \begin {scriptsize} Затверджено на засіданні кафедри теоретичної кібернетики, протокол \No 2  від   07.10.2024 р.\par\smallskip\smallskip\smallskip\smallskip
{\parbox {6.5 cm} {Зав. кафедри \dotfill Ю.В.~Крак}}\hspace {0.7cm} {\hbox to 6.5 cm { Екзаменатор \dotfill Т.О.~Карнаух}} \end{scriptsize}}
%end
\newpage
%begin
{{\begin {scriptsize}\par
{ \center  Київський національний університет імені Тараса Шевченка \par}
{\parbox {5 cm}{Освітня програма \hfill}{\parbox {7 cm}{Інформатика\hfill}}\parbox {6 cm}{\hfill Семестр 1 (маг)}}\par
{\parbox {5 cm} {Навчальний предмет \hfill}\parbox {7 cm}{Інформаційні мережі \hfill}\parbox {6cm}{\hfill}\par}\vspace{-15pt} \end{scriptsize}}{\center Екзаменаційний білет \No \textbf { 7 }\par}}
{
\qquad\textbf{Дати розгорнуту та вичерпну відповідь на питання}, висвітливши необхідні означення, властивості та за необхідності провівши порівняльний аналіз.

%beginitem
1. Ідентифікація, аутентифікація, авторизація. Технології аутентифікації, їх переваги та недоліки, а також способи реалізації.

%enditem


%\vspace{2mm}
\qquad\textbf{Визначити істинність таких тверджень}.

%beginitem
2. Модель OSI стосується мережі комутації пакетів.

%enditem

%beginitem
3. На фізичному, канальному та мережевому рівнях об'єкти взаємодіють через інтерфейси, а на транспортному та прикладному рівнях через протоколи.

%enditem

%beginitem
4. Протокол TCP можна використовувати для широкомовлення.

%enditem

%beginitem
5. Трафік відеоконференції є еластичним.

%enditem

%beginitem
6. Викладач дозволив на іспиті студентам використовувати тільки власні рукописні конспекти. Це можна віднести до авторизації.

%enditem
}
{\smallskip\smallskip \begin {scriptsize} Затверджено на засіданні кафедри теоретичної кібернетики, протокол \No 2  від   07.10.2024 р.\par\smallskip\smallskip\smallskip\smallskip
{\parbox {6.5 cm} {Зав. кафедри \dotfill Ю.В.~Крак}}\hspace {0.7cm} {\hbox to 6.5 cm { Екзаменатор \dotfill Т.О.~Карнаух}} \end{scriptsize}}
%end
\vspace{0.9cm}
%begin
{{\begin {scriptsize}\par
{ \center  Київський національний університет імені Тараса Шевченка \par}
{\parbox {5 cm}{Освітня програма \hfill}{\parbox {7 cm}{Інформатика\hfill}}\parbox {6 cm}{\hfill Семестр 1 (маг)}}\par
{\parbox {5 cm} {Навчальний предмет \hfill}\parbox {7 cm}{Інформаційні мережі \hfill}\parbox {6cm}{\hfill}\par}\vspace{-15pt} \end{scriptsize}}{\center Екзаменаційний білет \No \textbf { 8 }\par}}
{
\qquad\textbf{Дати розгорнуту та вичерпну відповідь на питання}, висвітливши необхідні означення, властивості та за необхідності провівши порівняльний аналіз.

%beginitem
1. Модель OSI. Рівні та їх задачі. Принципи будови моделі.

%enditem


%\vspace{2mm}
\qquad\textbf{Визначити істинність таких тверджень}.

%beginitem
2. Маршрутизатор зазвичай працює на мережевому рівні.

%enditem

%beginitem
3. Мережевий протокол містить інформацію про кількість даних, передану мережею за певну одиницю часу.

%enditem

%beginitem
4. Протокол SMTP має механізм контролю перенавантаження (congestion).

%enditem

%beginitem
5. Інтенсивність трафіка, рівна 1, вже є критичною для функціонування мережі.

%enditem

%beginitem
6. Використання допомоги товариша на іспиті є неавторизованим доступом до інформації.

%enditem
}
{\smallskip\smallskip \begin {scriptsize} Затверджено на засіданні кафедри теоретичної кібернетики, протокол \No 2  від   07.10.2024 р.\par\smallskip\smallskip\smallskip\smallskip
{\parbox {6.5 cm} {Зав. кафедри \dotfill Ю.В.~Крак}}\hspace {0.7cm} {\hbox to 6.5 cm { Екзаменатор \dotfill Т.О.~Карнаух}} \end{scriptsize}}
%end
\vspace{0.9cm}
%begin
{{\begin {scriptsize}\par
{ \center  Київський національний університет імені Тараса Шевченка \par}
{\parbox {5 cm}{Освітня програма \hfill}{\parbox {7 cm}{Інформатика\hfill}}\parbox {6 cm}{\hfill Семестр 1 (маг)}}\par
{\parbox {5 cm} {Навчальний предмет \hfill}\parbox {7 cm}{Інформаційні мережі \hfill}\parbox {6cm}{\hfill}\par}\vspace{-15pt} \end{scriptsize}}{\center Екзаменаційний білет \No \textbf { 9 }\par}}
{
\qquad\textbf{Дати розгорнуту та вичерпну відповідь на питання}, висвітливши необхідні означення, властивості та за необхідності провівши порівняльний аналіз.

%beginitem
1. Характеристики інфокомунікаційних мереж.

%enditem


%\vspace{2mm}
\qquad\textbf{Визначити істинність таких тверджень}.

%beginitem
2. Маршрутизатор зазвичай працює на прикладному рівні.

%enditem

%beginitem
3. Мережевий протокол містить інформацію про відхилення параметрів мережі від заданих.

%enditem

%beginitem
4. У протоколі HTTP встановлюється з'єднання між прикладними процесами.

%enditem

%beginitem
5. Профілювання трафіку полягає в тім, що згладжується пульсація трафіку.

%enditem

%beginitem
6. В аудиторію зайшов співробітник ректорату і представився. Це він виконав ідентифікацію.

%enditem
}
{\smallskip\smallskip \begin {scriptsize} Затверджено на засіданні кафедри теоретичної кібернетики, протокол \No 2  від   07.10.2024 р.\par\smallskip\smallskip\smallskip\smallskip
{\parbox {6.5 cm} {Зав. кафедри \dotfill Ю.В.~Крак}}\hspace {0.7cm} {\hbox to 6.5 cm { Екзаменатор \dotfill Т.О.~Карнаух}} \end{scriptsize}}
%end
\newpage
%begin
{{\begin {scriptsize}\par
{ \center  Київський національний університет імені Тараса Шевченка \par}
{\parbox {5 cm}{Освітня програма \hfill}{\parbox {7 cm}{Інформатика\hfill}}\parbox {6 cm}{\hfill Семестр 1 (маг)}}\par
{\parbox {5 cm} {Навчальний предмет \hfill}\parbox {7 cm}{Інформаційні мережі \hfill}\parbox {6cm}{\hfill}\par}\vspace{-15pt} \end{scriptsize}}{\center Екзаменаційний білет \No \textbf { 10 }\par}}
{
\qquad\textbf{Дати розгорнуту та вичерпну відповідь на питання}, висвітливши необхідні означення, властивості та за необхідності провівши порівняльний аналіз.

%beginitem
1. Стек протоколів TCP/IP. Співвідношення з моделлю OSI.

%enditem


%\vspace{2mm}
\qquad\textbf{Визначити істинність таких тверджень}.

%beginitem
2. Маршрутизатор зазвичай працює на канальному (data link) рівні.

%enditem

%beginitem
3. Мережевий протокол є фізичним інтерфейсом між об'єктами одного рівня.

%enditem

%beginitem
4. Протокол TCP підтримує багатоадресну передачу даних.

%enditem

%beginitem
5. Коефіцієнт пульсації знаходиться в межах від 0 (не включаючи) до 1.

%enditem

%beginitem
6. Викладач дозволив на іспиті студентам використовувати тільки власні рукописні конспекти. Це можна віднести до аутентифікації.

%enditem
}
{\smallskip\smallskip \begin {scriptsize} Затверджено на засіданні кафедри теоретичної кібернетики, протокол \No 2  від   07.10.2024 р.\par\smallskip\smallskip\smallskip\smallskip
{\parbox {6.5 cm} {Зав. кафедри \dotfill Ю.В.~Крак}}\hspace {0.7cm} {\hbox to 6.5 cm { Екзаменатор \dotfill Т.О.~Карнаух}} \end{scriptsize}}
%end
\vspace{0.9cm}
%begin
{{\begin {scriptsize}\par
{ \center  Київський національний університет імені Тараса Шевченка \par}
{\parbox {5 cm}{Освітня програма \hfill}{\parbox {7 cm}{Інформатика\hfill}}\parbox {6 cm}{\hfill Семестр 1 (маг)}}\par
{\parbox {5 cm} {Навчальний предмет \hfill}\parbox {7 cm}{Інформаційні мережі \hfill}\parbox {6cm}{\hfill}\par}\vspace{-15pt} \end{scriptsize}}{\center Екзаменаційний білет \No \textbf { 11 }\par}}
{
\qquad\textbf{Дати розгорнуту та вичерпну відповідь на питання}, висвітливши необхідні означення, властивості та за необхідності провівши порівняльний аналіз.

%beginitem
1. Методи забезпечення належного рівня QoS. Методи боротьби з перенавантаженням.

%enditem


%\vspace{2mm}
\qquad\textbf{Визначити істинність таких тверджень}.

%beginitem
2. Маршрутизатор зазвичай працює на транспортному рівні.

%enditem

%beginitem
3. Мережевий протокол є логічним інтерфейсом між різними рівнями стеку протоколів.

%enditem

%beginitem
4. Протокол  HTTP є протоколом без збереження стану (stateless).

%enditem

%beginitem
5. Мультимедійний трафік завжди є чутливим до спотворень.

%enditem

%beginitem
6. Використання допомоги товариша на іспиті є авторизованим доступом до інформації.

%enditem
}
{\smallskip\smallskip \begin {scriptsize} Затверджено на засіданні кафедри теоретичної кібернетики, протокол \No 2  від   07.10.2024 р.\par\smallskip\smallskip\smallskip\smallskip
{\parbox {6.5 cm} {Зав. кафедри \dotfill Ю.В.~Крак}}\hspace {0.7cm} {\hbox to 6.5 cm { Екзаменатор \dotfill Т.О.~Карнаух}} \end{scriptsize}}
%end
\vspace{0.9cm}
%begin
{{\begin {scriptsize}\par
{ \center  Київський національний університет імені Тараса Шевченка \par}
{\parbox {5 cm}{Освітня програма \hfill}{\parbox {7 cm}{Інформатика\hfill}}\parbox {6 cm}{\hfill Семестр 1 (маг)}}\par
{\parbox {5 cm} {Навчальний предмет \hfill}\parbox {7 cm}{Інформаційні мережі \hfill}\parbox {6cm}{\hfill}\par}\vspace{-15pt} \end{scriptsize}}{\center Екзаменаційний білет \No \textbf { 12 }\par}}
{
\qquad\textbf{Дати розгорнуту та вичерпну відповідь на питання}, висвітливши необхідні означення, властивості та за необхідності провівши порівняльний аналіз.

%beginitem
1. QoS. Основні типи QoS. Угода про якість обслуговування (SLA).

%enditem


%\vspace{2mm}
\qquad\textbf{Визначити істинність таких тверджень}.

%beginitem
2. Модель OSI стосується мережі комутації каналів.

%enditem

%beginitem
3. Мережевий протокол є логічним інтерфейсом між об'єктами одного рівня.

%enditem

%beginitem
4. При передачі файлів за протоколом TCP за втрати пакетів зазвичай виконується повторне відправлення втрачених пакетів.

%enditem

%beginitem
5. Пересічний користувач під час веб-сьорфінгу породжує потоковий трафік.

%enditem

%beginitem
6. В аудиторію зайшов співробітник ректорату і представився. Це він виконав авторизацію.

%enditem
}
{\smallskip\smallskip \begin {scriptsize} Затверджено на засіданні кафедри теоретичної кібернетики, протокол \No 2  від   07.10.2024 р.\par\smallskip\smallskip\smallskip\smallskip
{\parbox {6.5 cm} {Зав. кафедри \dotfill Ю.В.~Крак}}\hspace {0.7cm} {\hbox to 6.5 cm { Екзаменатор \dotfill Т.О.~Карнаух}} \end{scriptsize}}
%end
\newpage
%begin
{{\begin {scriptsize}\par
{ \center  Київський національний університет імені Тараса Шевченка \par}
{\parbox {5 cm}{Освітня програма \hfill}{\parbox {7 cm}{Інформатика\hfill}}\parbox {6 cm}{\hfill Семестр 1 (маг)}}\par
{\parbox {5 cm} {Навчальний предмет \hfill}\parbox {7 cm}{Інформаційні мережі \hfill}\parbox {6cm}{\hfill}\par}\vspace{-15pt} \end{scriptsize}}{\center Екзаменаційний білет \No \textbf { 13 }\par}}
{
\qquad\textbf{Дати розгорнуту та вичерпну відповідь на питання}, висвітливши необхідні означення, властивості та за необхідності провівши порівняльний аналіз.

%beginitem
1. Маршрутизація, її задачі. Алгоритми та протоколи маршрутизації.

%enditem


%\vspace{2mm}
\qquad\textbf{Визначити істинність таких тверджень}.

%beginitem
2. Маршрутизатор зазвичай працює на прикладному рівні.

%enditem

%beginitem
3. Мережевий протокол є логічним інтерфейсом між різними рівнями стеку протоколів.

%enditem

%beginitem
4. Естафетна передача виникає тоді, коли дані з однієї локальної мережі надходять в іншу.

%enditem

%beginitem
5. Засоби відеоконференцій зазвичай породжують потоковий трафік.

%enditem

%beginitem
6. Використання ChatGPT на іспиті є авторизованим доступом до інформації.

%enditem
}
{\smallskip\smallskip \begin {scriptsize} Затверджено на засіданні кафедри теоретичної кібернетики, протокол \No 2  від   07.10.2024 р.\par\smallskip\smallskip\smallskip\smallskip
{\parbox {6.5 cm} {Зав. кафедри \dotfill Ю.В.~Крак}}\hspace {0.7cm} {\hbox to 6.5 cm { Екзаменатор \dotfill Т.О.~Карнаух}} \end{scriptsize}}
%end
\vspace{0.9cm}
%begin
{{\begin {scriptsize}\par
{ \center  Київський національний університет імені Тараса Шевченка \par}
{\parbox {5 cm}{Освітня програма \hfill}{\parbox {7 cm}{Інформатика\hfill}}\parbox {6 cm}{\hfill Семестр 1 (маг)}}\par
{\parbox {5 cm} {Навчальний предмет \hfill}\parbox {7 cm}{Інформаційні мережі \hfill}\parbox {6cm}{\hfill}\par}\vspace{-15pt} \end{scriptsize}}{\center Екзаменаційний білет \No \textbf { 14 }\par}}
{
\qquad\textbf{Дати розгорнуту та вичерпну відповідь на питання}, висвітливши необхідні означення, властивості та за необхідності провівши порівняльний аналіз.

%beginitem
1. Ідентифікація, аутентифікація, авторизація. Технології аутентифікації, їх переваги та недоліки, а також способи реалізації.

%enditem


%\vspace{2mm}
\qquad\textbf{Визначити істинність таких тверджень}.

%beginitem
2. Маршрутизатор зазвичай працює на мережевому рівні.

%enditem

%beginitem
3. На фізичному, канальному та мережевому рівнях об'єкти взаємодіють через інтерфейси, а на транспортному та прикладному рівнях через протоколи.

%enditem

%beginitem
4. На прикладному рівні встановлюється з'єднання між мережевими інтерфейсами пристроїв.

%enditem

%beginitem
5. Протоколи, призначені для передачі мультимедіа, допускають втрату частини пакетів.

%enditem

%beginitem
6. В аудиторію зайшов співробітник ректорату і представився. Це він виконав ідентифікацію.

%enditem
}
{\smallskip\smallskip \begin {scriptsize} Затверджено на засіданні кафедри теоретичної кібернетики, протокол \No 2  від   07.10.2024 р.\par\smallskip\smallskip\smallskip\smallskip
{\parbox {6.5 cm} {Зав. кафедри \dotfill Ю.В.~Крак}}\hspace {0.7cm} {\hbox to 6.5 cm { Екзаменатор \dotfill Т.О.~Карнаух}} \end{scriptsize}}
%end
\vspace{0.9cm}
%begin
{{\begin {scriptsize}\par
{ \center  Київський національний університет імені Тараса Шевченка \par}
{\parbox {5 cm}{Освітня програма \hfill}{\parbox {7 cm}{Інформатика\hfill}}\parbox {6 cm}{\hfill Семестр 1 (маг)}}\par
{\parbox {5 cm} {Навчальний предмет \hfill}\parbox {7 cm}{Інформаційні мережі \hfill}\parbox {6cm}{\hfill}\par}\vspace{-15pt} \end{scriptsize}}{\center Екзаменаційний білет \No \textbf { 15 }\par}}
{
\qquad\textbf{Дати розгорнуту та вичерпну відповідь на питання}, висвітливши необхідні означення, властивості та за необхідності провівши порівняльний аналіз.

%beginitem
1. Стек протоколів TCP/IP. Співвідношення з моделлю OSI.

%enditem


%\vspace{2mm}
\qquad\textbf{Визначити істинність таких тверджень}.

%beginitem
2. Маршрутизатор зазвичай працює на транспортному рівні.

%enditem

%beginitem
3. Мережевий протокол містить інформацію про кількість даних, передану мережею за певну одиницю часу.

%enditem

%beginitem
4. На прикладному рівні встановлюється з'єднання між процесами.

%enditem

%beginitem
5. Засоби відеоконференцій зазвичай породжують пульсуючий трафік.

%enditem

%beginitem
6. Викладач дозволив на іспиті студентам використовувати тільки власні рукописні конспекти. Це можна віднести до авторизації.

%enditem
}
{\smallskip\smallskip \begin {scriptsize} Затверджено на засіданні кафедри теоретичної кібернетики, протокол \No 2  від   07.10.2024 р.\par\smallskip\smallskip\smallskip\smallskip
{\parbox {6.5 cm} {Зав. кафедри \dotfill Ю.В.~Крак}}\hspace {0.7cm} {\hbox to 6.5 cm { Екзаменатор \dotfill Т.О.~Карнаух}} \end{scriptsize}}
%end
\newpage
%begin
{{\begin {scriptsize}\par
{ \center  Київський національний університет імені Тараса Шевченка \par}
{\parbox {5 cm}{Освітня програма \hfill}{\parbox {7 cm}{Інформатика\hfill}}\parbox {6 cm}{\hfill Семестр 1 (маг)}}\par
{\parbox {5 cm} {Навчальний предмет \hfill}\parbox {7 cm}{Інформаційні мережі \hfill}\parbox {6cm}{\hfill}\par}\vspace{-15pt} \end{scriptsize}}{\center Екзаменаційний білет \No \textbf { 16 }\par}}
{
\qquad\textbf{Дати розгорнуту та вичерпну відповідь на питання}, висвітливши необхідні означення, властивості та за необхідності провівши порівняльний аналіз.

%beginitem
1. Модель OSI. Рівні та їх задачі. Принципи будови моделі.

%enditem


%\vspace{2mm}
\qquad\textbf{Визначити істинність таких тверджень}.

%beginitem
2. Маршрутизатор зазвичай працює на канальному (data link) рівні.

%enditem

%beginitem
3. Мережевий протокол є логічним інтерфейсом між об'єктами одного рівня.

%enditem

%beginitem
4. Використання протоколів із встановленням з'єднання можливе тільки за фізичного дротового прямого з'єднання хостів.

%enditem

%beginitem
5. Еластичний трафік не є чутливим до затримок.

%enditem

%beginitem
6. В аудиторію зайшов співробітник ректорату і представився. Це він виконав аутентифікацію.

%enditem
}
{\smallskip\smallskip \begin {scriptsize} Затверджено на засіданні кафедри теоретичної кібернетики, протокол \No 2  від   07.10.2024 р.\par\smallskip\smallskip\smallskip\smallskip
{\parbox {6.5 cm} {Зав. кафедри \dotfill Ю.В.~Крак}}\hspace {0.7cm} {\hbox to 6.5 cm { Екзаменатор \dotfill Т.О.~Карнаух}} \end{scriptsize}}
%end
\vspace{0.9cm}
%begin
{{\begin {scriptsize}\par
{ \center  Київський національний університет імені Тараса Шевченка \par}
{\parbox {5 cm}{Освітня програма \hfill}{\parbox {7 cm}{Інформатика\hfill}}\parbox {6 cm}{\hfill Семестр 1 (маг)}}\par
{\parbox {5 cm} {Навчальний предмет \hfill}\parbox {7 cm}{Інформаційні мережі \hfill}\parbox {6cm}{\hfill}\par}\vspace{-15pt} \end{scriptsize}}{\center Екзаменаційний білет \No \textbf { 17 }\par}}
{
\qquad\textbf{Дати розгорнуту та вичерпну відповідь на питання}, висвітливши необхідні означення, властивості та за необхідності провівши порівняльний аналіз.

%beginitem
1. Маршрутизація, її задачі. Алгоритми та протоколи маршрутизації.

%enditem


%\vspace{2mm}
\qquad\textbf{Визначити істинність таких тверджень}.

%beginitem
2. Модель OSI стосується мережі комутації пакетів.

%enditem

%beginitem
3. Мережевий протокол є фізичним інтерфейсом між об'єктами одного рівня.

%enditem

%beginitem
4. Протокол UDP має механізм контролю перенавантаження (congestion).

%enditem

%beginitem
5. Інтенсивність трафіка, рівна 1, вже є критичною для функціонування мережі.

%enditem

%beginitem
6. Співробітник ректорату на іспиті вирішив перевірити документи студентів. Це він авторизує студентів.

%enditem
}
{\smallskip\smallskip \begin {scriptsize} Затверджено на засіданні кафедри теоретичної кібернетики, протокол \No 2  від   07.10.2024 р.\par\smallskip\smallskip\smallskip\smallskip
{\parbox {6.5 cm} {Зав. кафедри \dotfill Ю.В.~Крак}}\hspace {0.7cm} {\hbox to 6.5 cm { Екзаменатор \dotfill Т.О.~Карнаух}} \end{scriptsize}}
%end
\vspace{0.9cm}
%begin
{{\begin {scriptsize}\par
{ \center  Київський національний університет імені Тараса Шевченка \par}
{\parbox {5 cm}{Освітня програма \hfill}{\parbox {7 cm}{Інформатика\hfill}}\parbox {6 cm}{\hfill Семестр 1 (маг)}}\par
{\parbox {5 cm} {Навчальний предмет \hfill}\parbox {7 cm}{Інформаційні мережі \hfill}\parbox {6cm}{\hfill}\par}\vspace{-15pt} \end{scriptsize}}{\center Екзаменаційний білет \No \textbf { 18 }\par}}
{
\qquad\textbf{Дати розгорнуту та вичерпну відповідь на питання}, висвітливши необхідні означення, властивості та за необхідності провівши порівняльний аналіз.

%beginitem
1. Ідентифікація, аутентифікація, авторизація. Технології аутентифікації, їх переваги та недоліки, а також способи реалізації.

%enditem


%\vspace{2mm}
\qquad\textbf{Визначити істинність таких тверджень}.

%beginitem
2. Модель OSI стосується мережі комутації каналів.

%enditem

%beginitem
3. Мережевий протокол містить інформацію про відхилення параметрів мережі від заданих.

%enditem

%beginitem
4. При передачі файлів за протоколом UDP за втрати пакетів зазвичай виконується повторне відправлення втрачених пакетів.

%enditem

%beginitem
5. Протоколи, призначені для передачі мультимедіа, допускають втрату частини пакетів.

%enditem

%beginitem
6. Використання ChatGPT на іспиті є неавторизованим доступом до інформації.

%enditem
}
{\smallskip\smallskip \begin {scriptsize} Затверджено на засіданні кафедри теоретичної кібернетики, протокол \No 2  від   07.10.2024 р.\par\smallskip\smallskip\smallskip\smallskip
{\parbox {6.5 cm} {Зав. кафедри \dotfill Ю.В.~Крак}}\hspace {0.7cm} {\hbox to 6.5 cm { Екзаменатор \dotfill Т.О.~Карнаух}} \end{scriptsize}}
%end
\newpage
%begin
{{\begin {scriptsize}\par
{ \center  Київський національний університет імені Тараса Шевченка \par}
{\parbox {5 cm}{Освітня програма \hfill}{\parbox {7 cm}{Інформатика\hfill}}\parbox {6 cm}{\hfill Семестр 1 (маг)}}\par
{\parbox {5 cm} {Навчальний предмет \hfill}\parbox {7 cm}{Інформаційні мережі \hfill}\parbox {6cm}{\hfill}\par}\vspace{-15pt} \end{scriptsize}}{\center Екзаменаційний білет \No \textbf { 19 }\par}}
{
\qquad\textbf{Дати розгорнуту та вичерпну відповідь на питання}, висвітливши необхідні означення, властивості та за необхідності провівши порівняльний аналіз.

%beginitem
1. QoS. Основні типи QoS. Угода про якість обслуговування (SLA).

%enditem


%\vspace{2mm}
\qquad\textbf{Визначити істинність таких тверджень}.

%beginitem
2. Маршрутизатор зазвичай працює на транспортному рівні.

%enditem

%beginitem
3. Мережевий протокол є фізичним інтерфейсом між об'єктами одного рівня.

%enditem

%beginitem
4. Протокол UDP може працювати як із встановленням з'єднання, так і без встановлення з'єднання.

%enditem

%beginitem
5. Засоби відеоконференцій зазвичай породжують потоковий трафік.

%enditem

%beginitem
6. Співробітник ректорату на іспиті вирішив перевірити документи студентів. Це він автентифікує студентів.

%enditem
}
{\smallskip\smallskip \begin {scriptsize} Затверджено на засіданні кафедри теоретичної кібернетики, протокол \No 2  від   07.10.2024 р.\par\smallskip\smallskip\smallskip\smallskip
{\parbox {6.5 cm} {Зав. кафедри \dotfill Ю.В.~Крак}}\hspace {0.7cm} {\hbox to 6.5 cm { Екзаменатор \dotfill Т.О.~Карнаух}} \end{scriptsize}}
%end
\vspace{0.9cm}
%begin
{{\begin {scriptsize}\par
{ \center  Київський національний університет імені Тараса Шевченка \par}
{\parbox {5 cm}{Освітня програма \hfill}{\parbox {7 cm}{Інформатика\hfill}}\parbox {6 cm}{\hfill Семестр 1 (маг)}}\par
{\parbox {5 cm} {Навчальний предмет \hfill}\parbox {7 cm}{Інформаційні мережі \hfill}\parbox {6cm}{\hfill}\par}\vspace{-15pt} \end{scriptsize}}{\center Екзаменаційний білет \No \textbf { 20 }\par}}
{
\qquad\textbf{Дати розгорнуту та вичерпну відповідь на питання}, висвітливши необхідні означення, властивості та за необхідності провівши порівняльний аналіз.

%beginitem
1. Методи забезпечення належного рівня QoS. Методи боротьби з перенавантаженням.

%enditem


%\vspace{2mm}
\qquad\textbf{Визначити істинність таких тверджень}.

%beginitem
2. Маршрутизатор зазвичай працює на канальному (data link) рівні.

%enditem

%beginitem
3. Мережевий протокол є логічним інтерфейсом між різними рівнями стеку протоколів.

%enditem

%beginitem
4. Протокол  HTTP допускає постійне з'єднання.

%enditem

%beginitem
5. Пересічний користувач під час веб-сьорфінгу породжує пульсуючий трафік.

%enditem

%beginitem
6. Співробітник ректорату на іспиті вирішив перевірити документи студентів. Це він ідентифікує студентів.

%enditem
}
{\smallskip\smallskip \begin {scriptsize} Затверджено на засіданні кафедри теоретичної кібернетики, протокол \No 2  від   07.10.2024 р.\par\smallskip\smallskip\smallskip\smallskip
{\parbox {6.5 cm} {Зав. кафедри \dotfill Ю.В.~Крак}}\hspace {0.7cm} {\hbox to 6.5 cm { Екзаменатор \dotfill Т.О.~Карнаух}} \end{scriptsize}}
%end
\vspace{0.9cm}
%begin
{{\begin {scriptsize}\par
{ \center  Київський національний університет імені Тараса Шевченка \par}
{\parbox {5 cm}{Освітня програма \hfill}{\parbox {7 cm}{Інформатика\hfill}}\parbox {6 cm}{\hfill Семестр 1 (маг)}}\par
{\parbox {5 cm} {Навчальний предмет \hfill}\parbox {7 cm}{Інформаційні мережі \hfill}\parbox {6cm}{\hfill}\par}\vspace{-15pt} \end{scriptsize}}{\center Екзаменаційний білет \No \textbf { 21 }\par}}
{
\qquad\textbf{Дати розгорнуту та вичерпну відповідь на питання}, висвітливши необхідні означення, властивості та за необхідності провівши порівняльний аналіз.

%beginitem
1. Характеристики інфокомунікаційних мереж.

%enditem


%\vspace{2mm}
\qquad\textbf{Визначити істинність таких тверджень}.

%beginitem
2. Модель OSI стосується мережі комутації каналів.

%enditem

%beginitem
3. Мережевий протокол є логічним інтерфейсом між об'єктами одного рівня.

%enditem

%beginitem
4. Протокол TCP забезпечує надійну передачу даних.

%enditem

%beginitem
5. Велика варіація затримки є критичною для передачі мультимедіа в режимі реального часу.

%enditem

%beginitem
6. В аудиторію зайшов співробітник ректорату і представився. Це він виконав авторизацію.

%enditem
}
{\smallskip\smallskip \begin {scriptsize} Затверджено на засіданні кафедри теоретичної кібернетики, протокол \No 2  від   07.10.2024 р.\par\smallskip\smallskip\smallskip\smallskip
{\parbox {6.5 cm} {Зав. кафедри \dotfill Ю.В.~Крак}}\hspace {0.7cm} {\hbox to 6.5 cm { Екзаменатор \dotfill Т.О.~Карнаух}} \end{scriptsize}}
%end
\newpage
%begin
{{\begin {scriptsize}\par
{ \center  Київський національний університет імені Тараса Шевченка \par}
{\parbox {5 cm}{Освітня програма \hfill}{\parbox {7 cm}{Інформатика\hfill}}\parbox {6 cm}{\hfill Семестр 1 (маг)}}\par
{\parbox {5 cm} {Навчальний предмет \hfill}\parbox {7 cm}{Інформаційні мережі \hfill}\parbox {6cm}{\hfill}\par}\vspace{-15pt} \end{scriptsize}}{\center Екзаменаційний білет \No \textbf { 22 }\par}}
{
\qquad\textbf{Дати розгорнуту та вичерпну відповідь на питання}, висвітливши необхідні означення, властивості та за необхідності провівши порівняльний аналіз.

%beginitem
1. QoS. Основні типи QoS. Угода про якість обслуговування (SLA).

%enditem


%\vspace{2mm}
\qquad\textbf{Визначити істинність таких тверджень}.

%beginitem
2. Маршрутизатор зазвичай працює на мережевому рівні.

%enditem

%beginitem
3. Мережевий протокол містить інформацію про відхилення параметрів мережі від заданих.

%enditem

%beginitem
4. Протокол UDP працює із встановлення з'єднання.

%enditem

%beginitem
5. За інтенсивності трафіка 2 мережа ще може нормально функціонувати, а за інтенсивності 5 та більше варто вживати термінових заходів для підтримки працездатності мережі.

%enditem

%beginitem
6. Використання допомоги товариша на іспиті є неавторизованим доступом до інформації.

%enditem
}
{\smallskip\smallskip \begin {scriptsize} Затверджено на засіданні кафедри теоретичної кібернетики, протокол \No 2  від   07.10.2024 р.\par\smallskip\smallskip\smallskip\smallskip
{\parbox {6.5 cm} {Зав. кафедри \dotfill Ю.В.~Крак}}\hspace {0.7cm} {\hbox to 6.5 cm { Екзаменатор \dotfill Т.О.~Карнаух}} \end{scriptsize}}
%end
\vspace{0.9cm}
%begin
{{\begin {scriptsize}\par
{ \center  Київський національний університет імені Тараса Шевченка \par}
{\parbox {5 cm}{Освітня програма \hfill}{\parbox {7 cm}{Інформатика\hfill}}\parbox {6 cm}{\hfill Семестр 1 (маг)}}\par
{\parbox {5 cm} {Навчальний предмет \hfill}\parbox {7 cm}{Інформаційні мережі \hfill}\parbox {6cm}{\hfill}\par}\vspace{-15pt} \end{scriptsize}}{\center Екзаменаційний білет \No \textbf { 23 }\par}}
{
\qquad\textbf{Дати розгорнуту та вичерпну відповідь на питання}, висвітливши необхідні означення, властивості та за необхідності провівши порівняльний аналіз.

%beginitem
1. Методи забезпечення належного рівня QoS. Методи боротьби з перенавантаженням.

%enditem


%\vspace{2mm}
\qquad\textbf{Визначити істинність таких тверджень}.

%beginitem
2. Модель OSI стосується мережі комутації пакетів.

%enditem

%beginitem
3. На фізичному, канальному та мережевому рівнях об'єкти взаємодіють через інтерфейси, а на транспортному та прикладному рівнях через протоколи.

%enditem

%beginitem
4. При передачі файлів за втрати пакетів протокол IP виконує повторне відправлення втрачених пакетів.

%enditem

%beginitem
5. Для передачі документів мережею велика варіація затримки некритична.

%enditem

%beginitem
6. Викладач дозволив на іспиті студентам використовувати тільки власні рукописні конспекти. Це можна віднести до аутентифікації.

%enditem
}
{\smallskip\smallskip \begin {scriptsize} Затверджено на засіданні кафедри теоретичної кібернетики, протокол \No 2  від   07.10.2024 р.\par\smallskip\smallskip\smallskip\smallskip
{\parbox {6.5 cm} {Зав. кафедри \dotfill Ю.В.~Крак}}\hspace {0.7cm} {\hbox to 6.5 cm { Екзаменатор \dotfill Т.О.~Карнаух}} \end{scriptsize}}
%end
\vspace{0.9cm}
%begin
{{\begin {scriptsize}\par
{ \center  Київський національний університет імені Тараса Шевченка \par}
{\parbox {5 cm}{Освітня програма \hfill}{\parbox {7 cm}{Інформатика\hfill}}\parbox {6 cm}{\hfill Семестр 1 (маг)}}\par
{\parbox {5 cm} {Навчальний предмет \hfill}\parbox {7 cm}{Інформаційні мережі \hfill}\parbox {6cm}{\hfill}\par}\vspace{-15pt} \end{scriptsize}}{\center Екзаменаційний білет \No \textbf { 24 }\par}}
{
\qquad\textbf{Дати розгорнуту та вичерпну відповідь на питання}, висвітливши необхідні означення, властивості та за необхідності провівши порівняльний аналіз.

%beginitem
1. Ідентифікація, аутентифікація, авторизація. Технології аутентифікації, їх переваги та недоліки, а також способи реалізації.

%enditem


%\vspace{2mm}
\qquad\textbf{Визначити істинність таких тверджень}.

%beginitem
2. Маршрутизатор зазвичай працює на прикладному рівні.

%enditem

%beginitem
3. Мережевий протокол містить інформацію про кількість даних, передану мережею за певну одиницю часу.

%enditem

%beginitem
4. BitTorrent – це протокол транспортного рівня.

%enditem

%beginitem
5. Трафік відеоконференції є еластичним.

%enditem

%beginitem
6. Використання допомоги товариша на іспиті є авторизованим доступом до інформації.

%enditem
}
{\smallskip\smallskip \begin {scriptsize} Затверджено на засіданні кафедри теоретичної кібернетики, протокол \No 2  від   07.10.2024 р.\par\smallskip\smallskip\smallskip\smallskip
{\parbox {6.5 cm} {Зав. кафедри \dotfill Ю.В.~Крак}}\hspace {0.7cm} {\hbox to 6.5 cm { Екзаменатор \dotfill Т.О.~Карнаух}} \end{scriptsize}}
%end
\newpage
%begin
{{\begin {scriptsize}\par
{ \center  Київський національний університет імені Тараса Шевченка \par}
{\parbox {5 cm}{Освітня програма \hfill}{\parbox {7 cm}{Інформатика\hfill}}\parbox {6 cm}{\hfill Семестр 1 (маг)}}\par
{\parbox {5 cm} {Навчальний предмет \hfill}\parbox {7 cm}{Інформаційні мережі \hfill}\parbox {6cm}{\hfill}\par}\vspace{-15pt} \end{scriptsize}}{\center Екзаменаційний білет \No \textbf { 25 }\par}}
{
\qquad\textbf{Дати розгорнуту та вичерпну відповідь на питання}, висвітливши необхідні означення, властивості та за необхідності провівши порівняльний аналіз.

%beginitem
1. Характеристики інфокомунікаційних мереж.

%enditem


%\vspace{2mm}
\qquad\textbf{Визначити істинність таких тверджень}.

%beginitem
2. Модель OSI стосується мережі комутації пакетів.

%enditem

%beginitem
3. Мережевий протокол є логічним інтерфейсом між різними рівнями стеку протоколів.

%enditem

%beginitem
4. Протокол UDP забезпечує надійну передачу даних.

%enditem

%beginitem
5. Еластичний трафік є чутливим до затримок.

%enditem

%beginitem
6. Використання ChatGPT на іспиті є неавторизованим доступом до інформації.

%enditem
}
{\smallskip\smallskip \begin {scriptsize} Затверджено на засіданні кафедри теоретичної кібернетики, протокол \No 2  від   07.10.2024 р.\par\smallskip\smallskip\smallskip\smallskip
{\parbox {6.5 cm} {Зав. кафедри \dotfill Ю.В.~Крак}}\hspace {0.7cm} {\hbox to 6.5 cm { Екзаменатор \dotfill Т.О.~Карнаух}} \end{scriptsize}}
%end
\vspace{0.9cm}
%begin
{{\begin {scriptsize}\par
{ \center  Київський національний університет імені Тараса Шевченка \par}
{\parbox {5 cm}{Освітня програма \hfill}{\parbox {7 cm}{Інформатика\hfill}}\parbox {6 cm}{\hfill Семестр 1 (маг)}}\par
{\parbox {5 cm} {Навчальний предмет \hfill}\parbox {7 cm}{Інформаційні мережі \hfill}\parbox {6cm}{\hfill}\par}\vspace{-15pt} \end{scriptsize}}{\center Екзаменаційний білет \No \textbf { 26 }\par}}
{
\qquad\textbf{Дати розгорнуту та вичерпну відповідь на питання}, висвітливши необхідні означення, властивості та за необхідності провівши порівняльний аналіз.

%beginitem
1. Модель OSI. Рівні та їх задачі. Принципи будови моделі.

%enditem


%\vspace{2mm}
\qquad\textbf{Визначити істинність таких тверджень}.

%beginitem
2. Маршрутизатор зазвичай працює на канальному (data link) рівні.

%enditem

%beginitem
3. Мережевий протокол є логічним інтерфейсом між об'єктами одного рівня.

%enditem

%beginitem
4. З“єднання за протоколом  HTTP завжди є постійним.

%enditem

%beginitem
5. Еластичний трафік не є чутливим до затримок.

%enditem

%beginitem
6. В аудиторію зайшов співробітник ректорату і представився. Це він виконав аутентифікацію.

%enditem
}
{\smallskip\smallskip \begin {scriptsize} Затверджено на засіданні кафедри теоретичної кібернетики, протокол \No 2  від   07.10.2024 р.\par\smallskip\smallskip\smallskip\smallskip
{\parbox {6.5 cm} {Зав. кафедри \dotfill Ю.В.~Крак}}\hspace {0.7cm} {\hbox to 6.5 cm { Екзаменатор \dotfill Т.О.~Карнаух}} \end{scriptsize}}
%end
\vspace{0.9cm}
%begin
{{\begin {scriptsize}\par
{ \center  Київський національний університет імені Тараса Шевченка \par}
{\parbox {5 cm}{Освітня програма \hfill}{\parbox {7 cm}{Інформатика\hfill}}\parbox {6 cm}{\hfill Семестр 1 (маг)}}\par
{\parbox {5 cm} {Навчальний предмет \hfill}\parbox {7 cm}{Інформаційні мережі \hfill}\parbox {6cm}{\hfill}\par}\vspace{-15pt} \end{scriptsize}}{\center Екзаменаційний білет \No \textbf { 27 }\par}}
{
\qquad\textbf{Дати розгорнуту та вичерпну відповідь на питання}, висвітливши необхідні означення, властивості та за необхідності провівши порівняльний аналіз.

%beginitem
1. Стек протоколів TCP/IP. Співвідношення з моделлю OSI.

%enditem


%\vspace{2mm}
\qquad\textbf{Визначити істинність таких тверджень}.

%beginitem
2. Маршрутизатор зазвичай працює на прикладному рівні.

%enditem

%beginitem
3. Мережевий протокол містить інформацію про відхилення параметрів мережі від заданих.

%enditem

%beginitem
4. Протокол TCP має механізм контролю перенавантаження (congestion).

%enditem

%beginitem
5. Профілювання трафіку полягає в тім, що згладжується пульсація трафіку.

%enditem

%beginitem
6. Співробітник ректорату на іспиті вирішив перевірити документи студентів. Це він ідентифікує студентів.

%enditem
}
{\smallskip\smallskip \begin {scriptsize} Затверджено на засіданні кафедри теоретичної кібернетики, протокол \No 2  від   07.10.2024 р.\par\smallskip\smallskip\smallskip\smallskip
{\parbox {6.5 cm} {Зав. кафедри \dotfill Ю.В.~Крак}}\hspace {0.7cm} {\hbox to 6.5 cm { Екзаменатор \dotfill Т.О.~Карнаух}} \end{scriptsize}}
%end
\newpage
%begin
{{\begin {scriptsize}\par
{ \center  Київський національний університет імені Тараса Шевченка \par}
{\parbox {5 cm}{Освітня програма \hfill}{\parbox {7 cm}{Інформатика\hfill}}\parbox {6 cm}{\hfill Семестр 1 (маг)}}\par
{\parbox {5 cm} {Навчальний предмет \hfill}\parbox {7 cm}{Інформаційні мережі \hfill}\parbox {6cm}{\hfill}\par}\vspace{-15pt} \end{scriptsize}}{\center Екзаменаційний білет \No \textbf { 28 }\par}}
{
\qquad\textbf{Дати розгорнуту та вичерпну відповідь на питання}, висвітливши необхідні означення, властивості та за необхідності провівши порівняльний аналіз.

%beginitem
1. Маршрутизація, її задачі. Алгоритми та протоколи маршрутизації.

%enditem


%\vspace{2mm}
\qquad\textbf{Визначити істинність таких тверджень}.

%beginitem
2. Модель OSI стосується мережі комутації каналів.

%enditem

%beginitem
3. На фізичному, канальному та мережевому рівнях об'єкти взаємодіють через інтерфейси, а на транспортному та прикладному рівнях через протоколи.

%enditem

%beginitem
4. Усі протоколи транспортного рівня вимагають встановлення з'єднання.

%enditem

%beginitem
5. Передача документів породжує еластичний трафік.

%enditem

%beginitem
6. В аудиторію зайшов співробітник ректорату і представився. Це він виконав авторизацію.

%enditem
}
{\smallskip\smallskip \begin {scriptsize} Затверджено на засіданні кафедри теоретичної кібернетики, протокол \No 2  від   07.10.2024 р.\par\smallskip\smallskip\smallskip\smallskip
{\parbox {6.5 cm} {Зав. кафедри \dotfill Ю.В.~Крак}}\hspace {0.7cm} {\hbox to 6.5 cm { Екзаменатор \dotfill Т.О.~Карнаух}} \end{scriptsize}}
%end
\vspace{0.9cm}
%begin
{{\begin {scriptsize}\par
{ \center  Київський національний університет імені Тараса Шевченка \par}
{\parbox {5 cm}{Освітня програма \hfill}{\parbox {7 cm}{Інформатика\hfill}}\parbox {6 cm}{\hfill Семестр 1 (маг)}}\par
{\parbox {5 cm} {Навчальний предмет \hfill}\parbox {7 cm}{Інформаційні мережі \hfill}\parbox {6cm}{\hfill}\par}\vspace{-15pt} \end{scriptsize}}{\center Екзаменаційний білет \No \textbf { 29 }\par}}
{
\qquad\textbf{Дати розгорнуту та вичерпну відповідь на питання}, висвітливши необхідні означення, властивості та за необхідності провівши порівняльний аналіз.

%beginitem
1. Ідентифікація, аутентифікація, авторизація. Технології аутентифікації, їх переваги та недоліки, а також способи реалізації.

%enditem


%\vspace{2mm}
\qquad\textbf{Визначити істинність таких тверджень}.

%beginitem
2. Маршрутизатор зазвичай працює на мережевому рівні.

%enditem

%beginitem
3. Мережевий протокол є фізичним інтерфейсом між об'єктами одного рівня.

%enditem

%beginitem
4. З“єднання за протоколом  HTTP завжди є непостійним.

%enditem

%beginitem
5. Коефіцієнт пульсації знаходиться в межах від 0 (не включаючи) до 1.

%enditem

%beginitem
6. Співробітник ректорату на іспиті вирішив перевірити документи студентів. Це він автентифікує студентів.

%enditem
}
{\smallskip\smallskip \begin {scriptsize} Затверджено на засіданні кафедри теоретичної кібернетики, протокол \No 2  від   07.10.2024 р.\par\smallskip\smallskip\smallskip\smallskip
{\parbox {6.5 cm} {Зав. кафедри \dotfill Ю.В.~Крак}}\hspace {0.7cm} {\hbox to 6.5 cm { Екзаменатор \dotfill Т.О.~Карнаух}} \end{scriptsize}}
%end
\vspace{0.9cm}
%begin
{{\begin {scriptsize}\par
{ \center  Київський національний університет імені Тараса Шевченка \par}
{\parbox {5 cm}{Освітня програма \hfill}{\parbox {7 cm}{Інформатика\hfill}}\parbox {6 cm}{\hfill Семестр 1 (маг)}}\par
{\parbox {5 cm} {Навчальний предмет \hfill}\parbox {7 cm}{Інформаційні мережі \hfill}\parbox {6cm}{\hfill}\par}\vspace{-15pt} \end{scriptsize}}{\center Екзаменаційний білет \No \textbf { 30 }\par}}
{
\qquad\textbf{Дати розгорнуту та вичерпну відповідь на питання}, висвітливши необхідні означення, властивості та за необхідності провівши порівняльний аналіз.

%beginitem
1. Характеристики інфокомунікаційних мереж.

%enditem


%\vspace{2mm}
\qquad\textbf{Визначити істинність таких тверджень}.

%beginitem
2. Маршрутизатор зазвичай працює на транспортному рівні.

%enditem

%beginitem
3. Мережевий протокол містить інформацію про кількість даних, передану мережею за певну одиницю часу.

%enditem

%beginitem
4. Протокол IP забезпечує надійну передачу даних.

%enditem

%beginitem
5. Пересічний користувач під час веб-сьорфінгу породжує потоковий трафік.

%enditem

%beginitem
6. Співробітник ректорату на іспиті вирішив перевірити документи студентів. Це він авторизує студентів.

%enditem
}
{\smallskip\smallskip \begin {scriptsize} Затверджено на засіданні кафедри теоретичної кібернетики, протокол \No 2  від   07.10.2024 р.\par\smallskip\smallskip\smallskip\smallskip
{\parbox {6.5 cm} {Зав. кафедри \dotfill Ю.В.~Крак}}\hspace {0.7cm} {\hbox to 6.5 cm { Екзаменатор \dotfill Т.О.~Карнаух}} \end{scriptsize}}
%end
\newpage
%begin
{{\begin {scriptsize}\par
{ \center  Київський національний університет імені Тараса Шевченка \par}
{\parbox {5 cm}{Освітня програма \hfill}{\parbox {7 cm}{Інформатика\hfill}}\parbox {6 cm}{\hfill Семестр 1 (маг)}}\par
{\parbox {5 cm} {Навчальний предмет \hfill}\parbox {7 cm}{Інформаційні мережі \hfill}\parbox {6cm}{\hfill}\par}\vspace{-15pt} \end{scriptsize}}{\center Екзаменаційний білет \No \textbf { 31 }\par}}
{
\qquad\textbf{Дати розгорнуту та вичерпну відповідь на питання}, висвітливши необхідні означення, властивості та за необхідності провівши порівняльний аналіз.

%beginitem
1. Стек протоколів TCP/IP. Співвідношення з моделлю OSI.

%enditem


%\vspace{2mm}
\qquad\textbf{Визначити істинність таких тверджень}.

%beginitem
2. Модель OSI стосується мережі комутації пакетів.

%enditem

%beginitem
3. Мережевий протокол є логічним інтерфейсом між об'єктами одного рівня.

%enditem

%beginitem
4. Протокол  HTTP є протоколом із збереженням стану (stateful).

%enditem

%beginitem
5. Засоби відеоконференцій зазвичай породжують пульсуючий трафік.

%enditem

%beginitem
6. Використання допомоги товариша на іспиті є неавторизованим доступом до інформації.

%enditem
}
{\smallskip\smallskip \begin {scriptsize} Затверджено на засіданні кафедри теоретичної кібернетики, протокол \No 2  від   07.10.2024 р.\par\smallskip\smallskip\smallskip\smallskip
{\parbox {6.5 cm} {Зав. кафедри \dotfill Ю.В.~Крак}}\hspace {0.7cm} {\hbox to 6.5 cm { Екзаменатор \dotfill Т.О.~Карнаух}} \end{scriptsize}}
%end
\vspace{0.9cm}
%begin
{{\begin {scriptsize}\par
{ \center  Київський національний університет імені Тараса Шевченка \par}
{\parbox {5 cm}{Освітня програма \hfill}{\parbox {7 cm}{Інформатика\hfill}}\parbox {6 cm}{\hfill Семестр 1 (маг)}}\par
{\parbox {5 cm} {Навчальний предмет \hfill}\parbox {7 cm}{Інформаційні мережі \hfill}\parbox {6cm}{\hfill}\par}\vspace{-15pt} \end{scriptsize}}{\center Екзаменаційний білет \No \textbf { 32 }\par}}
{
\qquad\textbf{Дати розгорнуту та вичерпну відповідь на питання}, висвітливши необхідні означення, властивості та за необхідності провівши порівняльний аналіз.

%beginitem
1. Маршрутизація, її задачі. Алгоритми та протоколи маршрутизації.

%enditem


%\vspace{2mm}
\qquad\textbf{Визначити істинність таких тверджень}.

%beginitem
2. Маршрутизатор зазвичай працює на мережевому рівні.

%enditem

%beginitem
3. Мережевий протокол містить інформацію про кількість даних, передану мережею за певну одиницю часу.

%enditem

%beginitem
4. Bluetooth є одним з протоколів стеку TCP/IP.

%enditem

%beginitem
5. Мультимедійний трафік завжди є чутливим до спотворень.

%enditem

%beginitem
6. Використання ChatGPT на іспиті є авторизованим доступом до інформації.

%enditem
}
{\smallskip\smallskip \begin {scriptsize} Затверджено на засіданні кафедри теоретичної кібернетики, протокол \No 2  від   07.10.2024 р.\par\smallskip\smallskip\smallskip\smallskip
{\parbox {6.5 cm} {Зав. кафедри \dotfill Ю.В.~Крак}}\hspace {0.7cm} {\hbox to 6.5 cm { Екзаменатор \dotfill Т.О.~Карнаух}} \end{scriptsize}}
%end
\vspace{0.9cm}
%begin
{{\begin {scriptsize}\par
{ \center  Київський національний університет імені Тараса Шевченка \par}
{\parbox {5 cm}{Освітня програма \hfill}{\parbox {7 cm}{Інформатика\hfill}}\parbox {6 cm}{\hfill Семестр 1 (маг)}}\par
{\parbox {5 cm} {Навчальний предмет \hfill}\parbox {7 cm}{Інформаційні мережі \hfill}\parbox {6cm}{\hfill}\par}\vspace{-15pt} \end{scriptsize}}{\center Екзаменаційний білет \No \textbf { 33 }\par}}
{
\qquad\textbf{Дати розгорнуту та вичерпну відповідь на питання}, висвітливши необхідні означення, властивості та за необхідності провівши порівняльний аналіз.

%beginitem
1. Модель OSI. Рівні та їх задачі. Принципи будови моделі.

%enditem


%\vspace{2mm}
\qquad\textbf{Визначити істинність таких тверджень}.

%beginitem
2. Маршрутизатор зазвичай працює на транспортному рівні.

%enditem

%beginitem
3. Мережевий протокол є фізичним інтерфейсом між об'єктами одного рівня.

%enditem

%beginitem
4. Естафетна передача обробляє перехід телефона між сотами.

%enditem

%beginitem
5. Трафік відеоконференції є еластичним.

%enditem

%beginitem
6. Використання допомоги товариша на іспиті є авторизованим доступом до інформації.

%enditem
}
{\smallskip\smallskip \begin {scriptsize} Затверджено на засіданні кафедри теоретичної кібернетики, протокол \No 2  від   07.10.2024 р.\par\smallskip\smallskip\smallskip\smallskip
{\parbox {6.5 cm} {Зав. кафедри \dotfill Ю.В.~Крак}}\hspace {0.7cm} {\hbox to 6.5 cm { Екзаменатор \dotfill Т.О.~Карнаух}} \end{scriptsize}}
%end
\newpage
\end{document}
